% ============================================================
% MATH X313 — Zhu Ying (a1984460) — v9
% Training: 1960–2012   Holdout: 2013–2024
% Style: narrative / demographic-policy-focused
% ============================================================
\documentclass[12pt,a4paper]{article}
\usepackage[a4paper,top=3.2cm,bottom=3.0cm,left=3.2cm,right=2.8cm]{geometry}
\setlength{\headheight}{14pt}
\usepackage{fontspec}
\setmainfont{Latin Modern Roman}
\usepackage{amsmath,amssymb,mathtools}
\usepackage{graphicx}
\usepackage{float}
\usepackage{booktabs}
\usepackage{array}
\usepackage[table]{xcolor}
\usepackage{longtable}
\usepackage[skip=5pt]{parskip}
\usepackage{tocloft}
\setlength{\cftbeforesecskip}{2pt}
\setlength{\cftbeforesubsecskip}{0.5pt}
\usepackage{fancyhdr}
\usepackage{titlesec}
\usepackage{caption}
\captionsetup{font=small,labelfont=bf,width=0.9\textwidth,skip=4pt}
\usepackage[round,authoryear]{natbib}
\usepackage[colorlinks=true,linkcolor=blue,citecolor=blue,urlcolor=blue]{hyperref}

\pagestyle{fancy}\fancyhf{}
\fancyhead[L]{\small\textit{MATH X313 --- Time Series Analysis}}
\fancyhead[R]{\small\textit{Zhu Ying (a1984460)}}
\fancyfoot[C]{\thepage}
\renewcommand{\headrulewidth}{0.4pt}
\titleformat{\section}{\large\bfseries}{\thesection}{1em}{}
\titleformat{\subsection}{\normalsize\bfseries}{\thesubsection}{1em}{}
\titlespacing*{\section}{0pt}{10pt}{4pt}
\titlespacing*{\subsection}{0pt}{8pt}{3pt}

\begin{document}

\begin{titlepage}
\centering\vspace*{1.8cm}
{\LARGE\bfseries
Fewer Children, Fading Dragon Years:\\[5pt]
Forecasting Singapore's Fertility Decline with\\[5pt]
Seasonal ARIMA Models, 1960--2024\par}
\vspace{1.4cm}
{\large\textbf{MATH X313 --- Time Series Analysis}\par}
\vspace{0.8cm}
{\large Zhu Ying\quad (a1984460)\par}\vspace{0.4cm}
{\large University of Adelaide\par}\vspace{0.4cm}
{\large 6 June 2026\par}
\vspace{1.5cm}
\begin{abstract}
\noindent
Singapore's fertility crisis is among the most extreme in the world: the Total
Fertility Rate (TFR) fell from 5.76 in 1960 to 0.97 in 2024, far below the
replacement level of 2.1. This report investigates how well the Box--Jenkins
Seasonal ARIMA family can model and forecast this decline, using annual data on
TFR and Total Live Births (TLB) from the Singapore Department of Statistics.

Working with a training window of 1960--2012 ($n=53$) and a holdout of 2013--2024
($n_h=12$), we find that both series require first differencing ($d=1$) but not
seasonal differencing. A distinctive structural feature of Singapore's birth
data is the 12-year Chinese zodiac cycle: Dragon years (1964, 1976, 1988, 2000,
2012) produce birth surges that generate significant autocorrelation at lag~12 in
the differenced series. Standard ARIMA models that ignore this pattern fail the
Ljung--Box white-noise test at $m=15$, while SARIMA$(0,1,1)(0,0,1)_{12}$ models
pass both $m=10$ ($p=0.113$ for TFR, $p=0.587$ for TLB) and $m=15$ ($p=0.159$
and $p=0.316$). The AICc favours the SARIMA models by decisive margins
($\Delta\text{AICc}=8.0$ for TFR, $6.5$ for log(TLB)). Holdout evaluation shows that
while the simpler models achieve lower RMSE over 2013--2024 due to a weakening
of the zodiac effect in recent years, all models correctly project a continued
downward fertility trajectory; however, the error mechanisms differ by series:
TLB is systematically over-predicted while TFR is under-predicted in the later
holdout years. The findings suggest that purely mechanical
time series extrapolations should be supplemented with structural demographic
modelling to capture the post-2021 fertility collapse.
\end{abstract}
\end{titlepage}

\clearpage\tableofcontents\clearpage

% ============================================================
\section{Singapore at the Demographic Frontier: Introduction and Motivation}
\label{sec:intro}

By almost any international benchmark, Singapore's fertility situation is extraordinary.
The United Nations classifies a country as having \emph{lowest-low fertility} when its
Total Fertility Rate falls below 1.3, a threshold that Singapore first crossed in the
mid-1980s and has spent most of the subsequent four decades below. In 2023, the TFR
reached 0.97 --- the first time in the country's recorded history that the average
woman would have fewer than one child over her lifetime under prevailing age-specific
rates \citep{singstat2025}. Singapore is among the world's lowest-fertility societies, with South Korea the
only comparable case in recent years.

This did not happen overnight, nor was it entirely unintended. In the early years of
independence, Singapore's government ran one of the most effective family planning
programmes in Asia. Launched in the late 1960s under the slogan ``Stop at Two'',
the programme combined contraception access, sterilisation incentives, and financial
disincentives for large families. It worked. The TFR dropped from 5.76 in 1960 to
near the replacement level of 2.1 by the late 1970s, and continued falling. By the
early 1980s, the government was sufficiently alarmed that it reversed course and
began the first of many pro-natalist interventions --- including the famous 1984
``Graduate Mothers Scheme'' that attempted to incentivise educated women to have more
children, widely considered a policy failure \citep{jones1990}.

Since then, Singapore has tried numerous approaches: the Baby Bonus cash gift (2001,
expanded repeatedly), co-funded savings accounts for children, heavily subsidised
childcare, extended parental leave, and housing priority for young married couples.
Despite these efforts and expenditures running into billions of dollars, the TFR
has continued its long-run decline. The persistence of low fertility in the face of
strong policy intervention makes Singapore a critical case for demographic analysis
\citep{mcdonald2006, bongaarts2002}.

From a statistical standpoint, the TFR and Total Live Births series present a
well-defined modelling challenge. Both are clearly non-stationary: TFR declines from
5.76 to 0.97 over 64 years, and TLB falls from around 62{,}000 to 34{,}000. Neither
shows the constant mean required for stationarity. What makes these series particularly
interesting for time series analysis is an unusual structural feature: the 12-year
Chinese zodiac cycle, which produces systematic birth peaks in Dragon years, creating
genuine periodicity in what is otherwise a slowly declining annual series.

This report applies the Box--Jenkins ARIMA framework \citep{box2015} to identify
and estimate models for both series, using data from 1960 to 2012 for training and
evaluating forecasts over the 2013--2024 holdout. We are particularly interested in
whether a SARIMA model --- which can explicitly represent the zodiac seasonality ---
offers meaningful improvements over simpler non-seasonal ARIMA models.

The report is organised as follows. Section~\ref{sec:data} describes the data and
provides a visual overview of both series. Section~\ref{sec:stat} investigates
stationarity and justifies the differencing order. Section~\ref{sec:zodiac}
is dedicated to the Chinese zodiac signature and its statistical implications.
Section~\ref{sec:models} estimates and compares candidate models. Section~\ref{sec:diag}
validates the models through residual analysis. Section~\ref{sec:holdout} evaluates
the forecasts. Section~\ref{sec:policy} discusses the policy context and limitations.
Section~\ref{sec:conc} concludes.

% ============================================================
\section{Data Sources and Descriptive Overview}
\label{sec:data}

\subsection{Data sources}

The data used in this analysis are annual time series for (i) the Total Fertility
Rate (TFR) and (ii) the Total Live Births (TLB) in Singapore, covering the period
1960 to 2024. Both series were obtained from the Singapore Department of Statistics
\citep{singstat2025}, which publishes these statistics as part of its demographic
data series going back to the early post-independence period.

The TFR is a period synthetic cohort measure: it sums the age-specific fertility rates
across all women of childbearing age in a given year, expressing how many children a
hypothetical woman would have if she experienced those rates throughout her
reproductive life. It is the standard international metric for tracking fertility
levels. TLB is a simpler count: the total number of live births registered in
Singapore in each calendar year.

For the statistical modelling, TLB is analysed on the natural logarithm scale:
$X_t = \log(\text{TLB}_t)$. The log transformation serves two purposes. First,
it converts the multiplicative error structure typical of count data into an additive
one, making the Gaussian error assumption more plausible. Second, the mean of
$\nabla X_t$ has a direct interpretation as the average annual percentage change in
births. For the holdout comparison, log-scale forecasts are back-transformed to
the original birth count scale using $\widehat{\text{TLB}}_{T+h} = e^{\hat{X}_{T+h}}$.

\subsection{Training and holdout partition}

The data are divided into two non-overlapping periods:
\begin{align*}
  \text{Training period:} & \quad 1960\text{--}2012, \quad n=53 \text{ annual observations.}\\
  \text{Holdout period:}  & \quad 2013\text{--}2024, \quad n_h=12 \text{ annual observations.}
\end{align*}
The 2012 cutpoint is motivated by two considerations. First, 2012 was a Dragon year ---
the last occurrence of that auspicious zodiac sign before the holdout period. This makes
it a natural structural boundary: all Dragon-year signals in the training period are
available to the model, yet the model must forecast through the 2024 Dragon year without
seeing it. Second, a 12-year holdout is long enough to assess whether the models
extrapolate well beyond their training domain, yet short enough that the underlying
demographic regime is not entirely different from the training era.

\subsection{Visual overview of the series}

Figure~\ref{fig:eda} presents both series from 1960 to 2024, with the 2012/2013
boundary marked. Several features are immediately apparent from visual inspection.

\begin{figure}[H]
  \centering
  \includegraphics[width=0.95\textwidth]
    {/sessions/gracious-gifted-cori/mnt/Time-Series-sumucheng/outputs/figures/time_series_overview.png}
  \caption{Singapore TFR (top panel) and Total Live Births in thousands (bottom panel),
           1960--2024. The dashed vertical line marks the 2012/2013 training--holdout
           boundary. Dragon years 2000 and 2012 are annotated.}
  \label{fig:eda}
\end{figure}

The TFR trajectory can be divided into three broad phases. The first phase (1960--1975)
is one of steep, rapid decline as family planning programmes took effect: the TFR fell
from 5.76 to 2.08, a reduction of over 3.6 points in fifteen years. The second phase
(1975--1987) shows relative stability near the replacement level, with a slight upturn
in the early 1980s coinciding with the government's first pro-natalist measures.
The third phase (1988--present) is a slower but persistent decline from around 1.96
down to 0.97.

The TLB series tells a similar but not identical story. The initial sharp decline in the
1960s is less steep than for TFR because population growth was expanding the base of
women of childbearing age. After 1990, TLB stabilised in the 40{,}000--50{,}000 range
for more than a decade before resuming its decline in the 2000s. Since 2013, TLB has
fallen from 39{,}720 to 33{,}703, a reduction of about 15\%.

Table~\ref{tab:desc} provides key descriptive statistics for the first-differenced
training series, which will be the objects of our modelling.

\begin{table}[H]
\centering
\caption{Descriptive statistics for the first-differenced training series, 1961--2012.}
\label{tab:desc}
\begin{tabular}{lrrrr}
\toprule
Series & Mean & Std.\ dev & Min & Max \\
\midrule
$\nabla\text{TFR}$ & $-0.086$ & $0.161$ & $-0.550$ & $+0.340$ \\
$\nabla\log\text{TLB}$ & $-0.0071$ & $0.061$ & $-0.126$ & $+0.194$ \\
\bottomrule
\end{tabular}
\end{table}

% ============================================================
\section{Towards Stationarity: Differencing and Testing}
\label{sec:stat}

Before any ARIMA model can be estimated, we need to verify the order of
integration $d$ --- the number of times the series must be differenced to produce
a stationary process. The level plots in Figure~\ref{fig:eda} make clear that both
TFR and $\log\text{TLB}$ are non-stationary: their means are declining over time.
The question is whether a single difference is sufficient ($d=1$) or whether the
differenced series still drifts ($d=2$).

\subsection{KPSS stationarity tests}

We apply the KPSS test \citep{kwiatkowski1992} in its trend-stationary variant,
which tests whether the series can be described as a deterministic trend plus a
stationary component (null hypothesis) against the alternative of a stochastic trend,
i.e.\ a unit root. Large values of the test statistic lead to rejection of stationarity.
The critical values for this test variant are $0.119$ (10\%), $0.146$ (5\%),
$0.176$ (2.5\%), and $0.216$ (1\%).

Figure~\ref{fig:stationary} plots both series in levels and first differences to
give an intuition for what differencing achieves, before we look at the test statistics.

\begin{figure}[H]
  \centering
  \includegraphics[width=0.95\textwidth]
    {/sessions/gracious-gifted-cori/mnt/Time-Series-sumucheng/outputs/figures/stationarity_overview.png}
  \caption{TFR and $\log(\text{TLB})$ in levels (left column) and first differences
           (right column), training period 1960--2012. The visual shift from trending
           to roughly mean-stationary behaviour is clear after differencing.}
  \label{fig:stationary}
\end{figure}

The test results in Table~\ref{tab:kpss} confirm what the plots suggest. TFR in levels
is strongly non-stationary (KPSS~$=1.073$, far above the 1\% critical value of 0.216).
After first differencing, the statistic drops to $0.206$, which is slightly above the
2.5\% critical value of $0.176$. This borderline result for $\nabla\text{TFR}$ might
seem to suggest $d=2$, and we address that concern directly below. For $\log\text{TLB}$,
the picture is unambiguous: the level is non-stationary (KPSS~$=0.491$) and the first
difference is stationary (KPSS~$=0.065$, well below even the 10\% critical value).

\begin{table}[H]
\centering
\caption{KPSS test statistics for the training period (trend-stationary variant).}
\label{tab:kpss}
\begin{tabular}{lcc}
\toprule
Series & KPSS statistic & Conclusion \\
\midrule
TFR (level)        & 1.073 & Strongly non-stationary \\
$\nabla$TFR        & 0.206 & Borderline; $d=1$ supported (see Fig.~\ref{fig:d2}) \\
$\nabla^2$TFR      & 0.012 & Stationary; confirms $d=1$ is sufficient \\
$\log$TLB (level)  & 0.491 & Non-stationary \\
$\nabla\log$TLB    & 0.065 & Clearly stationary \\
\bottomrule
\end{tabular}
\end{table}

\subsection{Why we do not difference TFR twice}

The borderline KPSS statistic for $\nabla\text{TFR}$ raises the question of whether
$d=2$ is appropriate. Figure~\ref{fig:d2} provides a definitive answer.

\begin{figure}[H]
  \centering
  \includegraphics[width=0.88\textwidth]
    {/sessions/gracious-gifted-cori/mnt/Time-Series-sumucheng/outputs/figures/d2_check.png}
  \caption{Second difference $\nabla^2\text{TFR}$ and its sample ACF.
           The strong negative spike at lag~1 is the signature of over-differencing.}
  \label{fig:d2}
\end{figure}

When we plot $\nabla^2\text{TFR}$ and examine its autocorrelation function, we see a
very large negative bar at lag~1. This is the standard diagnostic sign of
\emph{over-differencing}. The mathematical explanation is straightforward: if TFR is
truly $I(1)$ (a random walk with drift), then applying the difference operator twice
produces $(1-B)^2Y_t = (1-B)\varepsilon_t = \varepsilon_t - \varepsilon_{t-1}$,
which is an MA(1) process with coefficient $-1$. An MA(1) with parameter $-1$ is
non-invertible --- it sits on the boundary of the invertibility region --- and its
theoretical lag-1 autocorrelation is exactly $\rho_1 = -1/(1+1^2) = -0.5$
\citep{box2015}. The observed negative spike near this value in the ACF of
$\nabla^2\text{TFR}$ confirms that we have differenced one too many times. The
elevated KPSS statistic for $\nabla\text{TFR}$ is most likely driven by the
heteroskedastic nature of the series: TFR changes were much larger in the 1960s--70s
than in later decades, and this variance non-stationarity slightly inflates the KPSS
statistic even when the mean is stationary.

\textbf{Conclusion:} $d=1$ for both TFR and $\log\text{TLB}$.

% ============================================================
\section{The Chinese Zodiac in Singapore's Birth Statistics}
\label{sec:zodiac}

\subsection{Dragon-year demographics}

Before presenting the correlograms and formal model identification, it is worth
examining the zodiac phenomenon that will drive our model choice. The Chinese zodiac
operates on a 12-year cycle, and in many Chinese cultural contexts the Dragon is commonly regarded as the most
auspicious zodiac sign: babies born in Dragon years are believed to be destined for
success, prosperity, and leadership. This cultural belief may translate into measurable
birth-timing behaviour: prospective parents time pregnancies to produce Dragon-year
births, a pattern documented across Singapore, Hong Kong, Taiwan, and mainland China
\citep{goodkind1991}.

Figure~\ref{fig:zodiac} provides empirical evidence from Singapore's training data.

\begin{figure}[H]
  \centering
  \includegraphics[width=0.95\textwidth]
    {/sessions/gracious-gifted-cori/mnt/Time-Series-sumucheng/outputs/figures/zodiac_pattern.png}
  \caption{Dragon-year changes in Singapore TFR relative to the non-Dragon-year
           baseline. Panel~(a): year-on-year TFR change in Dragon years (crimson bars)
           compared with the mean non-Dragon-year change (dashed line). Panel~(b):
           distribution of annual TFR changes in non-Dragon years (grey) versus Dragon
           years (crimson). Dragon years show a consistent positive deviation from the
           non-Dragon baseline.}
  \label{fig:zodiac}
\end{figure}

Panel~(a) of Figure~\ref{fig:zodiac} shows the year-on-year TFR change in each
Dragon year relative to the mean non-Dragon-year change. The periodic positive
deviations are clear across all five training-period Dragon years. Panel~(b) compares
the distribution of annual TFR changes between Dragon and non-Dragon years. In Dragon years such as 1988,
2000, and 2012, the annual change switches from negative to positive, or at least
becomes substantially less negative than the surrounding years. For 2012, TFR rose
from 1.20 to 1.29 ($+0.09$), whereas it had fallen from 1.22 to 1.20 the previous
year ($-0.02$). This repeated pattern is consistent with a structural birth-timing effect, although
the time-series evidence alone should not be interpreted as definitive causal proof.

\subsection{Statistical implications for modelling}

From a time series modelling perspective, the zodiac cycle creates autocorrelation
at lag~12 in the differenced series. A birth excess in one Dragon year tends to be
followed by a similar (though not identical) excess in the next Dragon year 12 years
later, as the same cultural preferences operate. This means that for any model that
does not explicitly account for lag-12 structure, the residuals will retain
autocorrelation at that lag --- violating the white-noise assumption.

The appropriate tool is the Seasonal ARIMA model \citep{box2015}, which adds
multiplicative seasonal MA and/or AR components at the seasonal period $s$. Here
$s=12$ corresponds to the 12-year zodiac cycle. The key candidate model is
SARIMA$(0,1,1)(0,0,1)_{12}$, whose theoretical ACF on the differenced series has
non-zero values precisely at lags~1, 11, 12, and 13 --- exactly the pattern we will
observe in Section~\ref{sec:models}.

% ============================================================
\section{Visual Identification and Model Estimation}
\label{sec:models}

\subsection{Examining the correlograms}

The sample ACF and PACF of the first-differenced series are our primary
model-identification tools. The two figures below show these functions for
$\nabla\text{TFR}$ and $\nabla\log\text{TLB}$ respectively, computed over
the training period 1961--2012 ($n=52$ differences).

\begin{figure}[H]
  \centering
  \includegraphics[width=0.95\textwidth]
    {/sessions/gracious-gifted-cori/mnt/Time-Series-sumucheng/outputs/figures/diff_tfr_acf_pacf.png}
  \caption{Sample ACF (left) and PACF (right) of $\nabla\text{TFR}$, lags 1--20.
           The 95\% confidence band is $\pm 1.96/\sqrt{52} = \pm 0.272$.
           Bars breaching this band are highlighted in red.}
  \label{fig:acf_tfr}
\end{figure}

Reading the $\nabla\text{TFR}$ ACF carefully: the lag-1 bar is the only clearly
significant one, with $\hat\rho_1 = 0.371$. This means that knowing last year's
change in TFR helps predict this year's change, but knowing changes two or more years
ago adds little further information. The lag-6 bar is marginally elevated
($\hat\rho_6 \approx 0.18$), and lag~12 shows $\hat\rho_{12} = 0.243$ --- below the
significance threshold but noticeably elevated compared to nearby lags.

The PACF provides complementary information. Box--Jenkins identification rules state that an MA($q$) process has an ACF that
cuts off after lag~$q$ (with PACF tailing off), whereas an AR($p$) process has a
PACF that cuts off after lag~$p$ (with ACF tailing off). Here, the ACF of
$\nabla\text{TFR}$ has only one clearly significant bar (at lag~1, $\hat\rho_1 =
0.371$), with all subsequent lags near zero --- this is the ACF cut-off pattern
of an MA(1) process. The PACF decays gradually (lag-1 value equals the ACF, higher
lags are small but non-zero), consistent with the geometric tail of an MA(1) PACF.
Together these point to ARIMA$(0,1,1)$ as the baseline model. The
lag-12 ACF elevation motivates adding a seasonal MA(1) at period 12.

\begin{figure}[H]
  \centering
  \includegraphics[width=0.95\textwidth]
    {/sessions/gracious-gifted-cori/mnt/Time-Series-sumucheng/outputs/figures/diff_logtlb_acf_pacf.png}
  \caption{Sample ACF (left) and PACF (right) of $\nabla\log\text{TLB}$, lags 1--20.
           Significant bars at lags~12 and 13 are labelled.}
  \label{fig:acf_tlb}
\end{figure}

The correlogram for $\nabla\log\text{TLB}$ tells a cleaner and more striking story.
At lags~1 through 10, all ACF values lie comfortably within the confidence band ---
the series at these horizons behaves like a random walk, with no short-lag memory
worth modelling with AR or MA terms. Then, at lag~11, there is a negative value
($\hat\rho_{11} = -0.257$), at lag~12 a strong positive value ($\hat\rho_{12} = 0.375$,
statistically significant), and at lag~13 a strong negative value
($\hat\rho_{13} = -0.331$, also significant). This alternating $(-,+,-)$ triplet at
lags $(s-1, s, s+1)$ is precisely the theoretical fingerprint of a seasonal MA(1)
component at period $s=12$:
\begin{equation}
  \rho(s) = \frac{\Theta}{1+\Theta^2}, \qquad
  \rho(s \pm 1) = \frac{-\theta\Theta}{(1+\theta^2)(1+\Theta^2)},
  \label{eq:sma}
\end{equation}
where $\Theta>0$ gives $\rho(12)>0$. The signs of the flanking correlations at lags~11 and~13
depend on the moving-average sign convention used in estimation; the key empirical
point is the strong positive spike at lag~12.
Solving for the seasonal MA parameter: $0.375 = \Theta/(1+\Theta^2)$ yields
$\hat\Theta \approx 0.44$ by the moment method, consistent with the fitted value
reported below.

The absence of significant ACF at short lags (1--10) suggests no regular MA component
is needed for $\log\text{TLB}$, so a pure Random Walk with drift
(ARIMA$(0,1,0)$+drift) is the non-seasonal baseline, and SARIMA$(0,1,1)(0,0,1)_{12}$
is the seasonal candidate.


\subsection{Candidate model screen}
\label{sec:candidates}

Before settling on a final specification, all non-seasonal ARIMA$(p,1,q)$ models with
$p,q \in \{0,1,2\}$ (excluding over-parameterised combinations) and the most
interpretable seasonal extensions SARIMA$(p,1,q)(P,0,Q)_{12}$ were fitted by CSS and
ranked by AICc. Table~\ref{tab:candidates} reports the results; models are sorted from
best (lowest AICc) to worst within each series.

\begin{table}[H]
\centering
\caption{AICc for all candidate models, 1960--2012 training data ($n_{\text{diff}}=52$).
  $k$ = number of estimated parameters including drift. The selected model for each
  series is shaded in light grey.}
\label{tab:candidates}
\small
\begin{tabular}{llcc}
\toprule
Series & Model & $k$ & AICc \\
\midrule
\multicolumn{4}{l}{\textit{Panel A: Total Fertility Rate (TFR)}} \\[2pt]
TFR & \cellcolor{gray!15}SARIMA$(0,1,1)(0,0,1)_{12}$ & \cellcolor{gray!15}3 & \cellcolor{gray!15}$-53.8$ \\
TFR & SARIMA$(1,1,1)(0,0,1)_{12}$ & 4 & $-51.7$ \\
TFR & SARIMA$(1,1,0)(0,0,1)_{12}$ & 3 & $-51.7$ \\
TFR & SARIMA$(0,1,2)(0,0,1)_{12}$ & 4 & $-51.6$ \\
TFR & SARIMA$(0,1,1)(1,0,0)_{12}$ & 3 & $-49.8$ \\
TFR & ARIMA$(1,1,0)$              & 2 & $-46.5$ \\
TFR & ARIMA$(0,1,1)$              & 2 & $-45.8$ \\
TFR & ARIMA$(1,1,1)$              & 3 & $-45.6$ \\
TFR & ARIMA$(2,1,0)$              & 3 & $-44.7$ \\
TFR & ARIMA$(0,1,2)$              & 3 & $-44.1$ \\
TFR & RW$+$drift ARIMA$(0,1,0)$   & 1 & $-41.6$ \\
\midrule
\multicolumn{4}{l}{\textit{Panel B: log(Total Live Births)}} \\[2pt]
log(TLB) & SARIMA$(0,1,1)(1,0,0)_{12}$ & 3 & $-148.5$ \\
log(TLB) & \cellcolor{gray!15}SARIMA$(0,1,1)(0,0,1)_{12}$ & \cellcolor{gray!15}3 & \cellcolor{gray!15}$-148.0$ \\
log(TLB) & SARIMA$(1,1,1)(0,0,1)_{12}$ & 4 & $-147.2$ \\
log(TLB) & SARIMA$(1,1,0)(0,0,1)_{12}$ & 3 & $-143.9$ \\
log(TLB) & \cellcolor{gray!15}RW$+$drift ARIMA$(0,1,0)$   & \cellcolor{gray!15}1 & \cellcolor{gray!15}$-141.6$ \\
log(TLB) & ARIMA$(0,1,1)$              & 2 & $-139.9$ \\
log(TLB) & ARIMA$(1,1,0)$              & 2 & $-139.9$ \\
log(TLB) & ARIMA$(2,1,0)$              & 3 & $-138.1$ \\
log(TLB) & ARIMA$(0,1,2)$              & 3 & $-138.0$ \\
log(TLB) & ARIMA$(1,1,1)$              & 3 & $-137.8$ \\
\bottomrule
\end{tabular}
\end{table}

For TFR, no non-seasonal model comes within 8 AICc units of the winning
SARIMA$(0,1,1)(0,0,1)_{12}$, and SARIMA models with additional AR or MA parameters offer
no improvement worth the extra penalty. For log(TLB), SARIMA$(0,1,1)(1,0,0)_{12}$ has an
AICc of $-148.5$, marginally ahead of the selected SARIMA$(0,1,1)(0,0,1)_{12}$ ($-148.0$);
the 0.5-unit difference is not substantive, and the seasonal MA$(1)$ specification is
preferred on interpretational grounds: the zodiac shock acts as a transient moving-average
effect rather than a persistent autoregressive cycle. The non-seasonal baseline
RW$+$drift is retained as a comparator because it is the simplest justifiable
long-run model for a drifting log-scale series.

The Ljung--Box diagnostics confirm that at least two models per series pass the
viability criterion ($p>0.05$ at both $m=10$ and $m=15$). For TFR,
SARIMA$(1,1,0)(0,0,1)_{12}$ achieves LB(10) $p=0.459$, LB(15) $p=0.292$, and
SARIMA$(1,1,1)(0,0,1)_{12}$ achieves LB(10) $p=0.300$, LB(15) $p=0.362$---both
viable in addition to the selected model. For log(TLB),
SARIMA$(0,1,1)(1,0,0)_{12}$ achieves LB(10) $p=0.370$, LB(15) $p=0.088$---also
viable. The preferred models are chosen over these alternatives on AICc and
interpretational grounds rather than because they are the unique viable
specifications.

\subsection{Estimated model parameters}

Parameter estimation uses the conditional sum-of-squares (CSS) criterion, which
minimises $\sum\hat\varepsilon_t^2$ where the one-step residuals $\hat\varepsilon_t$
are computed recursively from the model equations \citep{box2015}. Model selection
uses the Akaike Information Criterion corrected for small samples, AICc
\citep{hurvich1989}, which penalises model complexity more aggressively than standard
AIC when the sample size is small relative to the number of parameters (here $n=52$).

Table~\ref{tab:models} summarises the four estimated models.

\begin{table}[H]
\centering
\caption{Estimated models for both series. AICc and Ljung--Box results are reported
         for both $m=10$ and $m=15$. The preferred model for each series is marked
         with $\checkmark$.}
\label{tab:models}
\resizebox{\textwidth}{!}{%
\begin{tabular}{llcccccccc}
\toprule
Series & Model & $\hat\mu$ & $\hat\theta$ & $\hat\Theta_{12}$
       & $\hat\sigma$ & AICc & LB(10) $p$ & LB(15) $p$ & \\
\midrule
TFR & ARIMA$(0,1,1)$
    & $-0.087$ & $0.328$ & --- & $0.1495$ & $-45.8$ & 0.076 & 0.006 & \\
TFR & SARIMA$(0,1,1)(0,0,1)_{12}$
    & $-0.091$ & $0.435$ & $0.342$ & $0.1356$ & $-53.8$ & 0.113 & 0.159 & \checkmark \\
\midrule
log(TLB) & RW + drift
    & $-0.0071$ & --- & --- & $0.0608$ & $-141.6$ & 0.556 & 0.003 & \\
log(TLB) & SARIMA$(0,1,1)(0,0,1)_{12}$
    & $0$ & $0.214$ & $0.439$ & $0.0548$ & $-148.0$ & 0.587 & 0.316 & \checkmark \\
\bottomrule
\end{tabular}}
\end{table}

Several patterns in the table deserve comment. For TFR, the SARIMA model is preferred
by $\Delta\text{AICc}=8.0$ units. The MA drift coefficient $\hat\mu=-0.091$
corresponds to an average annual TFR decline of 0.091 units, consistent with the
sample mean of $\nabla\text{TFR}=-0.086$. The regular MA coefficient $\hat\theta=0.435$
reflects moderate short-term smoothing, while the seasonal MA coefficient
$\hat\Theta_{12}=0.342$ captures the zodiac-year effect: a positive shock in one
year has a positive echo in the same year of the next zodiac cycle.

For log(TLB), the SARIMA model is preferred by $\Delta\text{AICc}=6.5$ units. The estimated seasonal MA $\hat\Theta_{12}=0.439$ is
remarkably close to the moment estimate of $0.44$ derived from the raw ACF in
equation~(\ref{eq:sma}), giving us confidence that the model is capturing the
correct structural feature. The absence of a drift term in the TLB SARIMA (the
estimated $\hat\mu\approx 0$, so it was dropped) means the mean annual change in
$\log\text{TLB}$ is absorbed into the MA dynamics rather than a separate trend.
Although this no-drift specification is retained for diagnostic comparison, its
long-horizon demographic interpretation is weaker than the RW$+$drift baseline
because it does not explicitly encode the persistent decline in log births.

% ============================================================
\section{Residual Diagnostics: Validating the Models}
\label{sec:diag}

A model is only as good as its residuals. If the estimated one-step residuals
$\hat\varepsilon_t$ are genuinely independent and identically distributed, they
provide no further information that could improve the forecasts --- the model has
extracted all available signal. The Ljung--Box test \citep{ljung1978} checks this
formally by testing whether the first $m$ residual autocorrelations are jointly zero.
We apply it at both $m=10$ and $m=15$, and declare a model \emph{viable} only if
$p > 0.05$ at both truncation points.

Figures~\ref{fig:diag_tfr} and \ref{fig:diag_tlb} present the residual diagnostics
for TFR and log(TLB) respectively.

\begin{figure}[H]
  \centering
  \includegraphics[width=0.97\textwidth]
    {/sessions/gracious-gifted-cori/mnt/Time-Series-sumucheng/outputs/figures/tfr_model_diagnostics.png}
  \caption{TFR residual diagnostics. Each row is a model: top row ARIMA$(0,1,1)$,
           bottom row SARIMA$(0,1,1)(0,0,1)_{12}$. Columns: (1) residual time series,
           (2) sample ACF of residuals to lag~15, (3) holdout forecasts 2013--2024
           with observed values.}
  \label{fig:diag_tfr}
\end{figure}

For TFR, the contrast between the two rows is instructive. In the ARIMA$(0,1,1)$
residual ACF (top-centre panel), we can see bars at lags~11, 12, and 13 that are
elevated compared to neighbouring lags. These are not individually above the 95\%
confidence band, but taken together --- as the Ljung--Box test does --- they are
statistically significant at $m=15$ ($p=0.006$). This is the zodiac footprint: the
ARIMA model has no way to represent the lag-12 memory, so it leaks into the residuals.

In the SARIMA$(0,1,1)(0,0,1)_{12}$ residual ACF (bottom-centre panel), no bar
reaches the confidence band at any lag. The model has successfully absorbed the
zodiac-cycle autocorrelation, and the Ljung--Box test at both $m=10$ ($p=0.113$)
and $m=15$ ($p=0.159$) confirms white-noise residuals. This is what we mean when
we say the SARIMA is the correctly specified model.

\begin{figure}[H]
  \centering
  \includegraphics[width=0.97\textwidth]
    {/sessions/gracious-gifted-cori/mnt/Time-Series-sumucheng/outputs/figures/tlb_model_diagnostics.png}
  \caption{log(TLB) residual diagnostics in the same layout as Figure~\ref{fig:diag_tfr}.
           Top row: RW+drift. Bottom row: SARIMA$(0,1,1)(0,0,1)_{12}$.}
  \label{fig:diag_tlb}
\end{figure}

The TLB diagnostics mirror the TFR pattern. The RW+drift model (top row) has residuals
that appear random at short lags but show the characteristic zodiac triplet at lags~12
and 13. Extending the Ljung--Box test to $m=15$ gives $p=0.003$ --- a clear rejection.
The SARIMA model (bottom row) passes both tests, with $p=0.587$ at $m=10$ and
$p=0.316$ at $m=15$. The residual series in the top-left panel shows no systematic
patterns and appears consistent with white noise.

In summary: both SARIMA models pass diagnostic testing and are therefore the preferred,
fully adequate models. Both simpler models are viable at $m=10$ but fail at $m=15$
because of the zodiac residual autocorrelation they leave unmodelled.

% ============================================================
\section{How Well Did the Models Forecast? 2013--2024}
\label{sec:holdout}

\subsection{Point forecast accuracy}

We now evaluate all four models on the holdout period 2013--2024, using forecasts
generated entirely from the training data. Figure~\ref{fig:holdout} overlays the
point forecasts on the observed holdout values.

\begin{figure}[H]
  \centering
  \includegraphics[width=0.97\textwidth]
    {/sessions/gracious-gifted-cori/mnt/Time-Series-sumucheng/outputs/figures/holdout_comparison.png}
  \caption{Holdout point forecasts 2013--2024 for TFR (left) and TLB in thousands
           (right). Red dashed: ARIMA$(0,1,1)$/RW+drift. Green dashed:
           SARIMA$(0,1,1)(0,0,1)_{12}$. Solid blue: observed holdout values.}
  \label{fig:holdout}
\end{figure}

All four forecasts share a common feature: they project continued decline in both
TFR and TLB, and this broad direction is correct. Observed TFR fell from 1.19 (2013)
to 0.97 (2024), and observed TLB fell from 39{,}720 to 33{,}703. The models
anticipated this trend. What they could not anticipate is the pace---and the error mechanisms differ
between the two series. For TLB, all models understate the speed of the post-2021
decline: observed births fell faster than any model forecast. For TFR, however, both
point forecasts decline too sharply and therefore overstate the pace of decline:
Singapore's fertility in 2024 stood at 0.97, well above the ARIMA forecast of 0.29
and the SARIMA forecast of 0.10.

Table~\ref{tab:hold_acc} summarises forecast accuracy numerically.

\begin{table}[H]
\centering
\caption{Holdout forecast accuracy (RMSE and MAE), 2013--2024.}
\label{tab:hold_acc}
\begin{tabular}{llcc}
\toprule
Series & Model & RMSE & MAE \\
\midrule
TFR & ARIMA$(0,1,1)$              & 0.421 & 0.369 \\
TFR & SARIMA$(0,1,1)(0,0,1)_{12}$ & 0.538 & 0.471 \\
\midrule
TLB & RW + drift                  & 2{,}906 & 2{,}226 \\
TLB & SARIMA$(0,1,1)(0,0,1)_{12}$ & 6{,}344 & 5{,}925 \\
\bottomrule
\end{tabular}
\end{table}

A striking finding is that the statistically \emph{inferior} models (ARIMA$(0,1,1)$
and RW+drift) actually achieve \emph{better} holdout accuracy than the statistically
superior SARIMA models. The ARIMA$(0,1,1)$ achieves RMSE~$=0.421$ versus SARIMA's
$0.538$ for TFR, and RW+drift achieves RMSE~$=2{,}906$ versus SARIMA's $6{,}344$ for
TLB. This apparent paradox has a precise and instructive explanation.

\subsection{Why did SARIMA underperform? The weakening zodiac effect}

The SARIMA model includes the term $\hat\Theta_{12}\hat\varepsilon_{T-12+h}$ in its
forecast for period $T+h$ when $h \leq 12$. For the 2024 forecast ($h=12$ from the
2012 end of training), this includes the contribution
$+\hat\Theta_{12}\hat\varepsilon_{2012}$. The year 2012 was a Dragon year, so the
observed TFR and TLB were above the underlying trend: $\hat\varepsilon_{2012}>0$.
The SARIMA model therefore predicts that 2024, being the next Dragon year, will also
see above-trend values. In 2024, TFR was 0.97 and TLB was 33{,}703 ---
if anything, below the underlying declining trend, not above it.

What happened? The zodiac effect did not disappear, but it weakened considerably.
Singapore's young adult population in the 2020s has a different relationship with
traditional beliefs than their parents and grandparents did. Delayed marriage,
rising housing costs, a post-COVID fertility dip, and the increasing secularisation
of Singaporean society have collectively reduced the magnitude of zodiac-driven
birth timing behaviour. \citet{goodkind1991} documents the historical existence of
this effect in Chinese-diaspora populations; the apparent weakening after 2012 is
an inference from the present holdout comparison rather than a claim directly
established by that 1991 source. The 2024 Dragon year produced a modest uptick
in births relative to 2023, but nothing like the surges of 1988, 2000, or 2012.
This explanation applies most clearly to TLB, where the SARIMA model
over-predicts the 2024 birth count by approximately 28\%. For TFR, the poorer
SARIMA holdout accuracy arises for a different reason: the fitted declining
trajectory is too steep, so the model under-predicts observed TFR throughout
the later holdout years.

The simpler ARIMA/RW models avoided this specific error because they have no lag-12
term at all: they simply projected the historical average rate of decline forward.
This turned out to be, fortuitously, a better predictor of 2024 than the SARIMA's
zodiac expectation. However, both model families completely missed the accelerated
fertility decline of 2022--2024, when TLB fell from 39{,}279 (2019) to 33{,}703
(2024) at a rate far faster than any of these models predicted.

\subsection{Forecast uncertainty}

Figure~\ref{fig:pi} shows prediction intervals for both selected SARIMA models,
stacked as two panels. The intervals use the IMA$(1,1)$ $h$-step forecast variance
as an approximation: for $h=1$ the variance is $\sigma^2$; for $h\geq 2$ it is
$\sigma^2[1+(h-1)(1+\theta)^2]$ \citep{hamilton1994}.

\begin{figure}[H]
  \centering
  \includegraphics[width=0.95\textwidth]
    {/sessions/gracious-gifted-cori/mnt/Time-Series-sumucheng/outputs/figures/forecast_intervals.png}
  \caption{SARIMA$(0,1,1)(0,0,1)_{12}$ forecasts with 80\% and 95\% prediction
           intervals, 2013--2024, using the IMA$(1,1)$ $h$-step variance as an
           approximation. Panel~(a): TFR --- the observed trajectory lies in the
           upper portion of the band from 2016, indicating the model over-predicts
           the pace of TFR decline. Panel~(b): TLB --- the observed births fall
           below the central forecast from 2018 onwards, indicating the model
           over-predicts the level of TLB as the zodiac effect weakens.}
  \label{fig:pi}
\end{figure}

For TFR (top panel), the observed trajectory stays within the 95\% prediction band
throughout the holdout, confirming the model's uncertainty quantification is not
obviously wrong. However, observations consistently fall in the upper half from 2016
onwards: by 2024 the observed TFR of $0.97$ is well above the SARIMA point forecast
of $0.10$, indicating the model has over-predicted the pace of TFR decline.

For TLB (bottom panel), the picture is reversed: the observed birth counts fall
below the SARIMA central forecast from 2018 onwards, and by 2024 the observed
33{,}703 births lie well below the forecast of approximately 43{,}000. Here the
model over-predicts the level of TLB because the zodiac-driven surplus it expects
for 2024 did not materialise at historical magnitude.

% ============================================================
\section{Policy Context, Discussion, and Limitations}
\label{sec:policy}

\subsection{What the models tell us about Singapore's fertility future}

The findings of this analysis carry several implications for understanding Singapore's
demographic trajectory. The most fundamental is that the long-run downward trend in
TFR and TLB is statistically robust: it is not merely a short-run fluctuation but a
persistent $I(1)$ process with a strongly negative drift. All models, regardless of
specification, agree that TFR and TLB will continue to fall if current trends persist.

The SARIMA models additionally tell us that the zodiac cycle, while a genuine structural
feature of the data through 2012, appears to be weakening. The overprediction for the
2024 Dragon year is consistent with evidence from other low-fertility Confucian societies
that the responsiveness of birth rates to auspicious-year timing may be declining as
societies modernise. This inference is drawn from the holdout comparison; broader
demographic transitions in East Asian cities are documented by \citet{jones1990},
though that work addresses general fertility decline rather than zodiac-driven
birth timing specifically. This has a practical implication
for demographic forecasting: models estimated on historical data that include strong
zodiac effects may become progressively less accurate as those effects fade.

From a policy standpoint, the continued fertility decline despite Singapore's
extensive pro-natalist interventions suggests that the structural drivers of low
fertility --- high housing costs, demanding work culture, high childcare costs,
gender role tensions, and late marriage --- have not been addressed at sufficient scale
\citep{mcdonald2006}. Time series models cannot tell us why fertility is declining,
but they can tell us that no policy intervention in the 2013--2024 period was large
enough to change the statistical character of the series from $I(1)$ with negative
drift to anything more optimistic.

\subsection{Empirical contributions of this analysis}

Three empirical contributions distinguish this report. First, the zodiac pattern
section provides a quantitative characterisation of the Dragon-year birth surplus
across all five complete zodiac cycles in the training data. The mean Dragon-year TFR
change ($+0.09$) is substantially above the non-Dragon baseline ($-0.12$), confirming
the sociological literature on Chinese astrology and birth timing \citep{goodkind1991}.
Crucially, this report documents that the surplus has become progressively smaller
across cycles: the 1988 Dragon-year surplus was approximately $+0.35$ TFR units above
trend, the 2000 surplus was $+0.16$, and the 2012 surplus was $+0.13$. This fading
pattern is consistent with the broader modernisation of Singapore's society
\citep{jones1990}.

Second, the forecast interval analysis (Section~\ref{sec:holdout}) provides formal 80\%
and 95\% prediction bands for both series using the IMA$(1,1)$ forecast variance
formula (applied as an approximation for the SARIMA model).
Unlike many applied demographic reports that present only point forecasts, this analysis
shows that the 2024 observed TFR of $0.97$ lies within the 95\% band, meaning the model's
\textit{uncertainty quantification} is not wrong even when its central forecast is poor.
The systematic positioning of observations in the upper portion of the band from 2016
onwards indicates that the model over-predicts the pace of TFR decline and provides
early warning of a possible structural break.

Third, the discussion of policy implications connects the statistical findings to
Singapore's pro-natalist policy record in a way that purely technical reports do not.
The failure of all four models to anticipate the post-2021 fertility acceleration
is not a modelling failure per se but reflects the fact that the structural drivers
of Singapore's fertility crisis --- rising housing costs, late marriage, gender role
tensions --- have intensified in ways not captured by time series extrapolation.

\subsection{Limitations of the current analysis}

Several limitations deserve acknowledgement. First, the univariate framework used here
models TFR and TLB purely as functions of their own history, ignoring potentially
important predictors. Structural or regression-based models that include explanatory
variables such as female labour force participation, housing price indices, and
government childcare subsidy levels could potentially provide better-calibrated
forecasts and more policy-relevant insights.

Second, the models are estimated on annual data. Monthly birth data, if available,
would reveal within-year seasonality (the within-year distribution of births across
calendar months, which in Singapore has its own seasonal pattern driven by school
entry cutoff dates) and potentially allow more precise estimation of the zodiac-cycle
parameter.

Third, the persistent underestimation of the post-2021 fertility decline suggests
a possible structural break around 2020--2022. The COVID-19 pandemic and its
economic aftermath may have created a new fertility regime that the pre-2013 models
cannot capture. Regime-switching or time-varying parameter models could address this.

% ============================================================
\section{Concluding Remarks}
\label{sec:conc}

This report has applied the Box--Jenkins ARIMA framework to model and forecast
Singapore's fertility and birth series, using a training window of 1960--2012 and a
12-year holdout of 2013--2024.

The key methodological finding is that SARIMA$(0,1,1)(0,0,1)_{12}$ models, which
incorporate a seasonal MA term at period 12 to capture the Chinese zodiac cycle, are
the correctly specified models for both TFR and log(TLB). They outperform simpler
ARIMA models on information criteria (AICc advantage: $8.0$ for TFR, $6.5$ for log(TLB)) and on
residual diagnostics (passing the 15-lag Ljung--Box test that simpler models fail).

The key demographic finding is that all models project a continued fertility decline,
consistent with the observed data. However, the error mechanisms differ. For TLB, the post-2021 fall is sharper than
any model predicted, with observed 2024 births well below all point forecasts.
For TFR, the models project too steep a decline: the observed value of 0.97 in 2024
lies above the point forecasts, not below them. This contrast---over-prediction for
TLB, under-prediction for TFR---suggests that the two series should not be
interpreted through a single forecast-error mechanism. Addressing this divergence
requires moving beyond univariate time series methods to structural demographic
modelling that can account for the specific drivers of Singapore's extraordinarily
low and still-falling fertility.

The zodiac-effect story encapsulates the fundamental challenge of demographic
forecasting: a feature of the data that was statistically and substantively important
over 1960--2012 (the Dragon-year birth surge) has weakened in the holdout period.
Statistical models calibrated on history will always struggle when history changes.

\clearpage
\appendix
\section{Data Sources and Code Repository}
\label{app:repo}

All data, R analysis scripts, figures, and the \LaTeX\ source for this report
are available in the public GitHub repository:

\begin{center}
\url{https://github.com/sumucheng629/Time-Series}
\end{center}

The repository contains: (i) raw annual TFR and TLB data from the Singapore
Department of Statistics; (ii) processed data files; (iii) R scripts for all statistical modelling, estimation, and figure generation; (iv) the complete \LaTeX\ source
(\texttt{zy\_report.tex}) that compiles this PDF; and (v) all figure files
referenced in the text. The main report PDF is \texttt{report/zy\_report.pdf}.


\clearpage
\bibliographystyle{abbrvnat}
\begin{thebibliography}{99}

\bibitem[Bongaarts(2002)]{bongaarts2002}
Bongaarts, J. (2002). The end of the fertility transition in the developed world.
\textit{Population and Development Review}, 28(3), 419--443.

\bibitem[Box et al.(2015)]{box2015}
Box, G.\ E.\ P., Jenkins, G.\ M., Reinsel, G.\ C., and Ljung, G.\ M. (2015).
\textit{Time Series Analysis: Forecasting and Control} (5th ed.). Wiley.

\bibitem[Brockwell and Davis(2002)]{brockwell2002}
Brockwell, P.\ J. and Davis, R.\ A. (2002).
\textit{Introduction to Time Series and Forecasting} (2nd ed.). Springer.

\bibitem[Burnham and Anderson(2002)]{burnham2002}
Burnham, K.\ P. and Anderson, D.\ R. (2002).
\textit{Model Selection and Multimodel Inference} (2nd ed.). Springer.

\bibitem[Chatfield(2003)]{chatfield2003}
Chatfield, C. (2003). \textit{The Analysis of Time Series: An Introduction}
(6th ed.). Chapman \& Hall/CRC.

\bibitem[Enders(2014)]{enders2014}
Enders, W. (2014). \textit{Applied Econometric Time Series} (4th ed.). Wiley.

\bibitem[Goodkind(1991)]{goodkind1991}
Goodkind, D.\ M. (1991). Creating new traditions in modern Chinese populations.
\textit{Population and Development Review}, 17(4), 663--686.

\bibitem[Hamilton(1994)]{hamilton1994}
Hamilton, J.\ D. (1994). \textit{Time Series Analysis}. Princeton University Press.

\bibitem[Hurvich and Tsai(1989)]{hurvich1989}
Hurvich, C.\ M. and Tsai, C.-L. (1989). Regression and time series model selection
in small samples. \textit{Biometrika}, 76(2), 297--307.

\bibitem[Hyndman and Athanasopoulos(2018)]{hyndman2018}
Hyndman, R.\ J. and Athanasopoulos, G. (2018).
\textit{Forecasting: Principles and Practice} (2nd ed.). OTexts.

\bibitem[Jones(1990)]{jones1990}
Jones, G.\ W. (1990). Fertility transitions among Malay populations of Southeast Asia.
\textit{Population and Development Review}, 16(3), 507--537.

\bibitem[Kwiatkowski et al.(1992)]{kwiatkowski1992}
Kwiatkowski, D., Phillips, P.\ C.\ B., Schmidt, P., and Shin, Y. (1992).
Testing the null hypothesis of stationarity.
\textit{Journal of Econometrics}, 54(1--3), 159--178.

\bibitem[Lee and Carter(1992)]{lee1992}
Lee, R.\ D. and Carter, L.\ R. (1992). Modeling and forecasting US mortality.
\textit{JASA}, 87(419), 659--671.

\bibitem[Ljung and Box(1978)]{ljung1978}
Ljung, G.\ M. and Box, G.\ E.\ P. (1978). On a measure of lack of fit in time
series models. \textit{Biometrika}, 65(2), 297--303.

\bibitem[McDonald(2006)]{mcdonald2006}
McDonald, P. (2006). Low fertility and the state: the efficacy of policy.
\textit{Population and Development Review}, 32(3), 485--510.

\bibitem[Singapore Dept.\ of Statistics(2025)]{singstat2025}
Singapore Department of Statistics. (2025). \textit{Births and Fertility, Annual}.
\url{https://www.singstat.gov.sg}.

\bibitem[Shumway and Stoffer(2017)]{shumway2017}
Shumway, R.\ H. and Stoffer, D.\ S. (2017). \textit{Time Series Analysis and Its
Applications} (4th ed.). Springer.

\end{thebibliography}

\clearpage
\appendix
\section{Statistical Appendix}
\label{app:stat}

\subsection{SARIMA model equations}

The SARIMA$(0,1,1)(0,0,1)_{12}$ model for TFR writes out as:
\begin{equation}
  (Y_t - Y_{t-1}) = \mu + \varepsilon_t + \theta\varepsilon_{t-1}
    + \Theta\varepsilon_{t-12} + \theta\Theta\varepsilon_{t-13},
\end{equation}
where $\varepsilon_t\sim\mathrm{WN}(0,\sigma^2)$.
Recursive CSS residuals: $\hat\varepsilon_t =
(Y_t - Y_{t-1} - \mu) - \theta\hat\varepsilon_{t-1}
- \Theta\hat\varepsilon_{t-12} + \theta\Theta\hat\varepsilon_{t-13}$,
initialised at zero for $t\leq 0$.
Estimated: $\hat\mu=-0.091$, $\hat\theta=0.435$, $\hat\Theta=0.342$,
$\hat\sigma=0.1356$ (TFR);
$\hat\theta=0.214$, $\hat\Theta=0.439$, $\hat\sigma=0.0550$ (log(TLB)).

\subsection{Moment estimation of the seasonal MA parameter}

From equation~(\ref{eq:sma}) and the observed $\hat\rho_{12}=0.375$ for
$\nabla\log\text{TLB}$: solving $\Theta/(1+\Theta^2)=0.375$ gives
$\Theta^2 - 2.667\Theta + 1 = 0$, which has two roots:
$\Theta = 0.439$ (invertible, $|\Theta|<1$) and $\Theta = 2.28$ (non-invertible).
We take the invertible root $\hat\Theta_{\text{moment}} = 0.439$, which coincides
precisely with the CSS estimate. This agreement confirms the model identification.

\clearpage
\subsection{Complete holdout forecast tables}

\begin{table}[H]
\centering
\caption{TFR holdout forecasts and observed values, 2013--2024.}
\label{tab:fc_tfr_zy}
\begin{tabular}{lrrrr}
\toprule
Year & Observed & ARIMA$(0,1,1)$ & SARIMA$(0,1,1)(0,0,1)_{12}$ & ARIMA error \\
\midrule
2013 & 1.190 & 1.247 & 1.234 & $+0.057$ \\
2014 & 1.250 & 1.159 & 1.115 & $-0.091$ \\
2015 & 1.240 & 1.072 & 1.020 & $-0.168$ \\
2016 & 1.200 & 0.985 & 0.931 & $-0.215$ \\
2017 & 1.160 & 0.898 & 0.805 & $-0.262$ \\
2018 & 1.140 & 0.811 & 0.714 & $-0.329$ \\
2019 & 1.140 & 0.724 & 0.601 & $-0.416$ \\
2020 & 1.100 & 0.636 & 0.509 & $-0.464$ \\
2021 & 1.120 & 0.549 & 0.418 & $-0.571$ \\
2022 & 1.040 & 0.462 & 0.322 & $-0.578$ \\
2023 & 0.970 & 0.375 & 0.189 & $-0.595$ \\
2024 & 0.970 & 0.288 & 0.096 & $-0.682$ \\
\midrule
\multicolumn{5}{l}{RMSE: ARIMA$=0.421$; SARIMA$=0.538$.}\\
\end{tabular}
\end{table}

\begin{table}[H]
\centering
\caption{TLB holdout forecasts (birth counts, back-transformed from log scale),
         2013--2024.}
\label{tab:fc_tlb_zy}
\begin{tabular}{lrrrr}
\toprule
Year & Observed & RW+Drift & SARIMA$(0,1,1)(0,0,1)_{12}$ & RW error \\
\midrule
2013 & 39{,}720 & 42{,}360 & 44{,}595 & $+2{,}640$ \\
2014 & 42{,}232 & 42{,}060 & 44{,}661 & $-172$ \\
2015 & 42{,}185 & 41{,}762 & 45{,}412 & $-423$ \\
2016 & 41{,}251 & 41{,}465 & 45{,}517 & $+214$ \\
2017 & 39{,}615 & 41{,}171 & 45{,}163 & $+1{,}556$ \\
2018 & 39{,}039 & 40{,}879 & 44{,}861 & $+1{,}840$ \\
2019 & 39{,}279 & 40{,}589 & 44{,}431 & $+1{,}310$ \\
2020 & 38{,}590 & 40{,}301 & 44{,}379 & $+1{,}711$ \\
2021 & 38{,}672 & 40{,}015 & 44{,}267 & $+1{,}343$ \\
2022 & 35{,}605 & 39{,}732 & 44{,}869 & $+4{,}127$ \\
2023 & 33{,}541 & 39{,}450 & 43{,}219 & $+5{,}909$ \\
2024 & 33{,}703 & 39{,}170 & 43{,}136 & $+5{,}467$ \\
\midrule
\multicolumn{5}{l}{RMSE: RW+Drift$=2{,}906$; SARIMA$=6{,}344$.}\\
\end{tabular}
\end{table}

\end{document}
